\chapter{Formal Sudoku Rules}

\section{State}\label{sec:state}

\begin{defn}
    The \emph{Sudoku size} of an $n \times n$ Sudoku puzzle is a positive integer $n$.
\end{defn}

\begin{defn}
    A \emph{Sudoku state} $U := \{U_{i,j} \mid i, j \in [n]\}$.
    Each element $U_{i,j} \subseteq [n]$.
\end{defn}

\begin{defn}
    The \emph{Sudoku task} $T := \{T_{i,j} \mid i, j \in [n]\}$ is the initial state of an unsolved Sudoku puzzle.
    Each element $T_{i,j}$ represents what Sudoku numbers can be in the cell $(i, j)$ of a Sudoku puzzle.
\end{defn}

\begin{defn}
    A \emph{Sudoku solution} $S$ is a Sudoku state such that every element $S_i$ has exactly one Sudoku number.
    It represents solved Sudoku puzzle.
    Each element $S_{i, j}$ represents what Sudoku number is in the cell $(i, j)$ of a solved Sudoku puzzle.
\end{defn}

\begin{defn}
    A Sudoku solution $S$ is \emph{valid} for a Sudoku puzzle with a task $T$, if
    $S_{i,j} \subseteq T_{i,j}$ for every $(i, j)$, and
    all Sudoku rules (all Sudoku regions and the variant) are satisfied.
\end{defn}

\begin{defn}
    \emph{$V$ is a substate of $U$}: $U \succeq V := \forall i, j \in [n]: U_{i, j} \supseteq V_{i, j}$
\end{defn}

\begin{defn}
    \emph{$V$ is a proper substate of $U$}: $U \succ V := U \succeq V \land U \neq V$
\end{defn}

\begin{defn}
    \emph{The difference of $U$ and $V$}: $(U \setminus V)_{i, j} := U_{i, j} \setminus V_{i, j}$
\end{defn}

\begin{defn}
    \emph{The union of $U$ and $V$}: $(U \cup V)_{i, j} := U_{i, j} \cup V_{i, j}$
\end{defn}

\begin{defn}
    \emph{The union of all valid solutions} is $S_{ALL} := \bigcup S_i$, where
    $S_i$ are all the valid solutions for a Sudoku puzzle with a task $T$.
\end{defn}

\section{Rules}

The definition in Section~\ref{sec:state} only specifies that
a valid solution chooses one number for each cell from the task.
Sudoku rules further restrict which solutions can be a valid solution.

Let us have a Sudoku puzzle with size $n$, task $T$, and a solution $S$.
Let $S_{ALL}$ be the union of all valid solutions.

\subsection*{Regions}

A Sudoku puzzle has a finite set of Sudoku regions $R := {R_1, R_2, \dots, R_r}$,
where a region $R_i \in ([n] \times [n])^n$ is an array of $n$ coordinates of cells of the Sudoku puzzle.

Each region $R_i$ adds an extra requirement on a valid solution of the puzzle:
Cells in a region must contain each Sudoku number exactly once.
(TODO: highlight the text of the rule?)

\subsection*{Variants}

A variant adds extra requirements for a valid solution that cannot be encoded by regions.

A variant only needs to include one validation function $V(S)$,
which takes a puzzle's solution and returns true if the solution does not break any rule of the variant
and returns false it does.
If validation function returns false, the solution is not valid.

\subsection*{Edge variant}

An edge variant is a specific type of variant.
All variants implemented in the software are edge variants.

Edge variant is $E := (L, E, E_L, L_N, V)$.

$L := [l]$
is a set of labels.

$E \subseteq \{(a, b) \mid a, b \in [n] \times [n] \land
\text{the cell } a \text{ is adjacent on the left or the top side of the cell } b\}$
is the set of some (oriented) pairs of edge-adjacent cells.
If $(a, b) \in E$, then the cells $a$ and $b$ are label-adjacent.
(TODO: definition?)
$E_L : E \to L$
is an edge-label function that for a label-adjacent pair of cells returns a label.
$L_N : L \to 2^{[n] \times [n]}$
is a label-neighbourhood function that for a label $l$ returns
the set of all (oriented) pairs of Sudoku numbers allowed in a
label-adjacent pair of cells with the label $l$.

$V$ is the validation function.
\[
    V(S) := \bigwedge_{(a, b) \in E} (S_a, S_b) \in L_N(E_L((a, b)))
\]
That means $V(S)$ returns true, if and only if for each pair of label-adjacent cells $(a, b) \in E$:
The pair of Sudoku numbers $(S_a, S_b)$ in the solution $S$ is allowed by the label-neighbourhood $L_N$ function
for the label $l := E_L((a, b))$.
